\documentclass[12pt]{article}
\usepackage[utf8]{inputenc}
\usepackage[T2A]{fontenc}
\usepackage[russian]{babel}
\usepackage{amsmath}
\usepackage{amssymb}
\usepackage{geometry}
\geometry{margin=2.5cm}

\title{Дифференциальный эффект Джоуля-Томсона в модели Дитеричи}

\date{}

\begin{document}

\maketitle

\tableofcontents
\newpage

\section{Введение}
Эффект Джоуля-Томсона описывает изменение температуры газа при адиабатическом расширении через пористую перегородку. Дифференциальный эффект Джоуля-Томсона ($\mu_{JT}$) определяется как производная температуры по давлению при постоянной энтальпии:
$$
\mu_{JT} = \left(\frac{\partial T}{\partial P}\right)_H.
$$
Если $\mu_{JT} > 0$, газ охлаждается при расширении; если $\mu_{JT} < 0$ — нагревается. В данном документе рассмотрим зависимость $\mu_{JT}$ от параметров модели Дитеричи.

\section{Общая формула для $\mu_{JT}$}
Для произвольного уравнения состояния $\mu_{JT}$ выражается через термодинамические величины:
$$
\mu_{JT} = \frac{1}{C_P} \left[T \left(\frac{\partial V}{\partial T}\right)_P - V\right],
$$
где $C_P$ — теплоемкость при постоянном давлении, $V$ — молярный объем, $T$ — температура.

\section{Уравнение Дитеричи}
Уравнение Дитеричи для реального газа:
$$
P(V - b) = RT e^{-\frac{a}{RTV}},
$$
где:

    
- $a$ — параметр, учитывающий силы притяжения между молекулами,
    
- $b$ — собственный объем молекул.


\section{Вывод $\mu_{JT}$ для модели Дитеричи}
Чтобы найти $\mu_{JT}$, необходимо вычислить $\left(\frac{\partial V}{\partial T}\right)_P$ из уравнения Дитеричи.

\subsection{Логарифмическое преобразование}
Прологарифмируем уравнение Дитеричи:
$$
\ln P + \ln(V - b) = \ln R + \ln T - \frac{a}{RTV}.
$$

\subsection{Дифференцирование по $T$ при постоянном $P$}
Дифференцируем обе стороны по $T$, считая $P = \text{const}$:
$$
\frac{\partial}{\partial T} \left[\ln(V - b)\right]_P = \frac{1}{T} + \frac{a}{RT^2V} - \frac{a}{RTV^2} \left(\frac{\partial V}{\partial T}\right)_P.
$$

\subsection{Решение относительно $\left(\frac{\partial V}{\partial T}\right)_P$}
После алгебраических преобразований:
$$
\left(\frac{\partial V}{\partial T}\right)_P = \frac{V(V - b)}{T} \left[1 + \frac{a}{RTV} - \frac{a}{RTV^2}(V - b)\right]^{-1}.
$$

\subsection{Подстановка в формулу $\mu_{JT}$}
Подставляем $\left(\frac{\partial V}{\partial T}\right)_P$ в выражение для $\mu_{JT}$:
$$
\mu_{JT} = \frac{1}{C_P} \left[T \cdot \frac{V(V - b)}{T} \left(1 + \frac{a}{RTV} - \frac{a}{RTV^2}(V - b)\right)^{-1} - V\right].
$$

\subsection{Упрощение выражения}
Итоговая формула:
$$
\mu_{JT} = \frac{V(V - b)}{C_P} \left[ \left(1 + \frac{a}{RTV} - \frac{a}{RTV^2}(V - b)\right)^{-1} - \frac{1}{V - b} \right].
$$

\section{Анализ зависимости $\mu_{JT}$ от параметров модели}

    
- \textbf{Параметр $a$:} Увеличение $a$ усиливает отрицательный вклад в $\mu_{JT}$, что способствует охлаждению газа ($\mu_{JT} > 0$).
    
- \textbf{Параметр $b$:} Увеличение $b$ уменьшает объем $V - b$, что усиливает положительный вклад в $\mu_{JT}$.
    
- \textbf{Температура $T$:} При высоких температурах ($T \to \infty$) эффект Джоуля-Томсона стремится к нулю, так как $a/(RTV) \to 0$, и уравнение Дитеричи переходит в идеальный газ.


\section{Сравнение с моделью Ван-дер-Ваальса}
В модели Ван-дер-Ваальса $\mu_{JT}$ выражается как:
$$
\mu_{JT} = \frac{1}{C_P} \left( \frac{2a}{RT} - b \right).
$$
В модели Дитеричи $\mu_{JT}$ зависит от экспоненциальной поправки $e^{-a/(RTV)}$, что делает зависимость от $T$ и $V$ более сложной. Это позволяет точнее описывать эффект Джоуля-Томсона при умеренных и высоких давлениях, где силы притяжения играют существенную роль.

\section{Физическая интерпретация}

    
- \textbf{Область охлаждения ($\mu_{JT} > 0$):} Реализуется, когда преобладают силы притяжения ($a$ доминирует над $b$). В модели Дитеричи эта область шире, чем в модели Ван-дер-Ваальса, благодаря экспоненциальной зависимости.
    
- \textbf{Температура инверсии:} Определяется из условия $\mu_{JT} = 0$. Для модели Дитеричи она зависит от комбинации параметров $a, b, T, V$, что позволяет точнее описывать экспериментальные данные.


\section{Заключение}
Модель Дитеричи учитывает межмолекулярные силы притяжения через экспоненциальную поправку, что делает её более точной для описания дифференциального эффекта Джоуля-Томсона в условиях высоких давлений и умеренных температур. По сравнению с моделью Ван-дер-Ваальса, она лучше предсказывает ширину области охлаждения и температуру инверсии, особенно для газов с сильными силами притяжения. Однако математическая сложность ограничивает её применение в инженерных расчетах, где предпочтительны более простые модели.

\end{document}